// Film Grain — density-domain photographic grain synthesis
//
// Models real film grain through optical density space, matching the
// physical process of silver halide exposure and dye-cloud formation.
//
// Pipeline:
//   1. Convert linear RGB to optical density: D = -log10(I)
//   2. Generate spatially-shaped noise (FBM for dye-cloud clumping)
//   3. Apply luminance-dependent amplitude (grain peaks in midtones,
//      tapers in highlights and deep shadows)
//   4. Cross-channel correlation (blue channel gets heavier grain,
//      matching real film where blue-sensitive emulsion has larger halides)
//   5. Add grain in density space, convert back to linear
//
// This produces grain that naturally "bites" into shadows and rolls off
// in highlights — the hallmark of real film stock, impossible to achieve
// with simple additive noise in gamma space.
//
// Connect @image in scene-linear color space.

f$strength = 0.35;       // grain intensity (0.0-1.0)
f$size = 1.0;            // grain size (1.0 = 35mm, larger = coarser)
f$roughness = 0.6;       // 0 = smooth dye clouds, 1 = sharp/gritty
f$color_variance = 0.5;  // 0 = monochrome grain, 1 = full RGB variation
f$shadow_response = 0.8; // grain amplitude in deep shadows
f$mid_response = 1.0;    // grain amplitude in midtones
f$highlight_response = 0.2; // grain amplitude in highlights

// ── Convert to optical density ───────────────────────────────────
vec3 img = @image;
float lum = luma(img);

// Density = -log10(linear). Clamp to avoid log(0).
float density = -log10(max(lum, 0.0001));
// Normalize density to roughly [0, 1] for the amplitude curve.
// Typical scene density range is 0 (white) to ~3 (deep shadow).
float d_norm = clamp(density / 3.0, 0.0, 1.0);

// ── Luminance-dependent amplitude (cubic spline approximation) ───
// Three control points: shadow (d=0.8), midtone (d=0.4), highlight (d=0.1)
// Blended via smoothstep transitions.
float amp_shadow = $shadow_response * smoothstep(0.5, 0.9, d_norm);
float amp_mid = $mid_response * (1.0 - abs(d_norm - 0.4) * 2.5);
amp_mid = max(amp_mid, 0.0);
float amp_highlight = $highlight_response * smoothstep(0.3, 0.0, d_norm);
float amplitude = max(amp_shadow + amp_mid + amp_highlight, 0.05) * $strength;

// ── Spatially-shaped noise (dye-cloud clumping) ──────────────────
// Use FBM at multiple scales to simulate halide cluster structure.
// The 'roughness' parameter blends between smooth low-octave clouds
// and sharp high-frequency grain.
float scale = 800.0 / max($size, 0.1);
float sx = u * scale;
float sy = v * scale;
// Frame-dependent seed for temporal independence (each frame = new grain)
float seed = fi * 137.31;

// Fine grain (high frequency, sharp particles)
float g_fine = simplex(sx + seed, sy + seed * 0.73);
// Medium grain (dye-cloud clumps)
float g_med = simplex(sx * 0.4 + seed * 1.1, sy * 0.4 + seed * 0.51);
// Coarse structure
float g_coarse = simplex(sx * 0.15 + seed * 0.87, sy * 0.15 + seed * 1.3);

// Blend scales based on roughness: high roughness = more fine grain
float grain_luma = lerp(g_med, g_fine, $roughness) * 0.7
                 + g_coarse * 0.3;

// ── Cross-channel correlation ────────────────────────────────────
// Real film: CMY dye layers are physically overlapping, creating
// correlated but not identical grain. Blue-sensitive layer has the
// largest halides (most grain). Fully monochrome at color_variance=0.
float g_r = simplex(sx + seed + 41.7, sy + seed + 17.3);
float g_g = simplex(sx + seed + 83.1, sy + seed + 59.7);
float g_b = simplex(sx + seed + 127.9, sy + seed + 91.1);

// Blend: correlated (mono) ←→ independent (color)
float gr_r = lerp(grain_luma, g_r, $color_variance);
float gr_g = lerp(grain_luma, g_g, $color_variance);
float gr_b = lerp(grain_luma, g_b, $color_variance * 0.7);

// Blue channel gets ~30% more grain (larger halides)
gr_b = gr_b * 1.3;

// ── Apply grain in density space ─────────────────────────────────
// Add noise to each channel's density, then convert back to linear.
float eps = 0.0001;
float dr = -log10(max(img.r, eps)) + gr_r * amplitude;
float dg = -log10(max(img.g, eps)) + gr_g * amplitude;
float db = -log10(max(img.b, eps)) + gr_b * amplitude;

// Convert back to linear light: I = 10^(-D)
vec3 result = vec3(
    pow10(-dr),
    pow10(-dg),
    pow10(-db)
);

@OUT = vec3(max(result.r, 0.0), max(result.g, 0.0), max(result.b, 0.0));
